\subsection{2. 预备知识}\label{preliminaries}

本节简要概述效应与余效应系统——支撑我们工作的两大理论支柱。我们假定读者
熟悉基本的类型理论与范畴论；本节的目的是固定记号，并引入若干关键抽象，
第~3 节将把这些抽象操作化为运行时机制。

\subsection{2.1. 效应}\label{efects}

在简单类型 λ 演算（STLC）{[}20, 21{]} 中，类型判断
\(\Gamma \vdash t : T\) 表明项 \(t\) 在上下文 $\Gamma$ 下具有类型
\(T\)。效应系统对类型加以细化，用以描述计算可能产生的副作用，从而得到
如下形式的判断：

\[
\Gamma \vdash t : T _ { \mathrm { e f f e c t } }\tag{1}
\]

这里，结果类型被附注以效应代数中的一个元素，用以刻画计算可能产生的副作
用，从而支持对带状态计算进行组合式推理。这一方法起源于 Lucassen 和
Gifford {[}22{]}，他们引入了一种带种类的类型系统，区分类型、效应与区
域，以发现并行程序中的调度约束。

单子效应。Moggi {[}16{]} 首次通过单子对计算效应进行了范畴论建模；Wadler
{[}23{]} 使这一方法在 Haskell 中得到推广。范畴 \(\mathcal{C}\) 上的一个
单子 \(( T , \eta , \mu )\) 将一次带效应的计算封装为类型 \(T ( A )\)
的一个值，其中 \(\eta : A \to T ( A )\) 提升纯值，
\(\mu : T ( T ( A ) ) \to T ( A )\) 将嵌套计算排序化。经典的实例包括
Maybe 单子（表示部分性）、State 单子（表示可变状态）以及 IO 单子（表示
外部交互）。

代数效应。Plotkin 和 Power {[}17, 24{]} 证明了代数操作可以决定单子，由
此建立了一个使效应接口与其实现相解耦的框架。一个效应签名 $\Sigma$ 声明
一组操作（例如状态用的
\(( \mathrm { e . g . , g e t : } ( ) \to S , \mathrm { p u t : } S \to ( )\)）；
程序可以自由地调用操作，而不必承诺特定的解释。随后，Plotkin 和 Pretnar
{[}25{]} 引入了效应处理器，通过提供续延语义来解释操作：

\[
\mathrm { h a n d l e } e \mathrm { w i t h } \left\{ \mathrm { o p } ( v , \kappa ) \mapsto \dots \right\}\tag{2}
\]

处理器接收操作实参 \(v\) 与定界续延 \(\kappa\)，并可以对续延调用零次、
一次或多次，从而在统一的框架内支持异常、协程与非确定性 {[}26{]}。诸如
Koka {[}27, 28{]}、Eff {[}29{]} 与 OCaml 5 {[}30{]} 等语言都已采用代
数效应，只是在设计权衡上各有取舍。

\subsection{2.2. 余效应}\label{coefects}

与效应相对偶，余效应系统 {[}18, 31{]} 所丰富的是上下文而非类型，得到如
下形式的判断：

\[
\Gamma _ { \mathrm { c o e f f e c t } } \vdash t : T\tag{3}
\]

这里，上下文被附注以余效应代数中的一个元素，用以刻画计算对其环境的要求，
例如需要访问的资源、需要持有的权限或需要依赖的服务。效应刻画程序对世界
的影响，而余效应刻画世界对程序的约束。

余单子余效应。用余单子来结构化依赖上下文的计算这一思想，最早由 Uustalu
和 Vene {[}32{]} 提出，他们提出了对称（半）幺半范畴上的余单子，作为
\(\mathrm { M o g g i ^ { \prime } s }\) 针对效应的单子框架的对偶，用以
刻画数据流与属性求值等概念。Petricek 等人 {[}18{]} 在此基础上提出余效
应，作为对上下文依赖性的统一静态分析。一个余单子 \(( D , \varepsilon ,
\delta )\) 刻画依赖上下文的计算：\(\varepsilon : D ( A ) \to A\) 从上下
文中提取当前值，\(\delta : D ( A ) \to D ( D ( A ) )\) 复制上下文以供嵌
套访问。Environment 余单子 \(D ( X ) = E \times X\) 刻画对固定环境 \(E\)
的依赖；Stream 余单子 \(D ( X ) = \mathbb { N } \to X\) 刻画对时序数据
的依赖。

分级余效应。为了更细粒度的追踪，分级余效应系统使用一个预序半环
\(\mathcal{S} = ( S , \leq , + , \times , 0 , 1 )\) 作为余效应代数
{[}33{]}；这一规范后来被 Gaboardi 等人 {[}19{]} 与分级效应统一起来。
\(S\) 的元素被用于附注每个变量绑定，以量化其使用量：0 表示未使用，1 表
示线性使用，\(n\) 表示有界使用，\(\infty\) 表示无限制使用。半环运算以顺
次（\(( \times )\)）与并行（\(( + )\)）两种方式组合余效应，从而在统一的
代数框架 {[}37{]} 内实现精确的资源追踪、敏感性分析 {[}34{]} 与信息流控
制 {[}35, 36{]}。

\subsection{2.3. 与动态组合的关系}\label{relationship-to-dynamic-composability}

效应与余效应系统沿两个互补的方向组织关于计算的推理：效应刻画计算如何修
改其环境，而余效应刻画计算如何依赖其环境。这两个方向正对应第~1 节所识别
的动态组合的两个维度：

• 时间维组合性要求组件对共享环境所做的修改在卸载时是可逆的。相关的效应
是带状态的那一类效应，它们持久地变换该环境；撤销这样的变换要求该变换存
在逆。

• 空间维组合性要求组件间的依赖以响应式的方式被声明和管理。这类依赖正是
余效应所刻画的对象，而对它们的管理在于将每一个依赖对照环境所提供的内容
加以解析。

然而，经典的效应与余效应系统是静态的工具：效应在词法固定的作用域内被追
踪，并由编译期的处理器来消解；余效应的标注则依据执行前即确定的上下文进
行校验。相比之下，动态组合要求这些保证对在运行时到达和离开的组件仍然成
立，而此时的上下文在持续演化。没有固定的词法作用域能够框定一个在部署之
后才加载的插件；没有编译期的上下文能够预先预见由运行时配置所涌现的依赖。

这促使我们转变视角：与其用更多标注去扩展静态类型系统，不如将效应与余效
应的概念结构具体化，使得运行时能够直接对它们进行操作，从而以动态的方式
建立起这些系统在静态时所提供的保证。
